El potassi 40 ($^{40}$K) és un isòtop inestable. Es pot transformar en calci (Ca) mitjançant una desintegració $\beta^-$ o en argó (Ar) mitjançant una desintegració $\beta^+$. El nombre atòmic del calci és 20.

\begin{apartat}{1}
Escriviu les equacions nuclears que corresponen a aquests processos, incloent-hi els neutrins i els antineutrins.
\end{apartat}

\begin{apartat}{1}
També és possible que el potassi 40 capturi un electró de la seva escorça i emeti un fotó gamma de 1\,460~MeV. Calculeu la longitud d'ona i la freqüència d'aquests raigs gamma. Calculeu també la disminució de la massa de l'àtom de potassi 40 deguda a l'energia que s'endú el fotó.
\end{apartat}

\dades{%
  Constant de Planck, $h = 6,63\times 10^{-34}\ \mathrm{J\,s}$.\\
  $|e| = 1,6\times 10^{-19}\ \mathrm{C}$.\\
  Velocitat de la llum, $c = 3,00\times 10^{8}\ \mathrm{m\,s^{-1}}$.\\
  $1\ \mathrm{eV} = 1,60\times 10^{-19}\ \mathrm{J}$.}
